\section{Workers, memory and responsive interaction}
\subsection{Why more workers can be slower}
Loading a field can involve filesystem reads, decompression, copying, metadata
validation and reductions. Several workers can overlap independent work, but
they also compete for memory bandwidth, cache capacity, storage and decoder
scratch. The number of logical CPU cores is therefore not a universal best
loading count. It is also unrelated to how many GPU compute queues the solver
should use.

On a new native desktop profile, Studio can run a short automatic measurement
after the interface becomes idle. It creates a small temporary float32 NPY
array, compares cached reads and scanning at candidate worker counts, and removes
the temporary fixture afterward. Each choice uses a warmup and three equal-work
trials. Automatic selects the smallest count close to the best measured result,
avoiding unnecessary concurrency for nearly identical measured throughput.

The measurement is bounded and cancelable. It is available through
\textbf{Settings → Performance → Measure this computer}, and the automatic
measurement can be disabled. The saved profile is reused only when its CPU,
core-count and RAM identity matches. An explicitly selected loading count
remains the user's choice; the profile applies when the setting is Automatic.

This is a cached NPY read-and-scan measurement, not a sustained disk benchmark,
ZFP or netCDF performance campaign, graphics benchmark or solver scaling test.
Different working sets, formats and background applications can change the best
count. Rerun it while the computer is otherwise idle, then check the actual
dataset's responsiveness. The native memory budget can still reduce how many
reads run concurrently below the selected worker count.

\subsection{The I/O memory preference}
\textbf{I/O memory limit (MiB)} controls temporary native field-reading memory.
Use zero for an automatic RAM-based policy or a manual value in the range shown
by Settings. A MiB is $2^{20}$ bytes. The permitted manual range begins at
64 MiB; its upper bound depends on available RAM and is capped by the implementation.
The active limits panel shows the applied limit, current reserved bytes and
waiting reads.

Saving a changed limit applies it without leaving the plot. When concurrent
reads temporarily occupy the budget, another read waits for capacity. If one
request cannot fit even by itself, it reports the required size rather than
waiting indefinitely. Increasing the limit retries applicable budget-failed
requests for the same selected field; a user does not need to leave a volume
view and return merely to trigger another read. Reducing the limit allows
existing reads to finish while new admissions respect the lower limit.

This budget is neither total application RAM nor video memory. Raising it cannot
increase a graphics texture limit, remove a contour primitive limit or enlarge
Python's address space. It can also increase simultaneous memory pressure, so
the useful setting is the one that admits the desired work while leaving room
for the rest of the application and operating system.

\subsection{The different pools of memory}
Studio has several owners with different lifetimes. The resource ledger names
native reads, decoded fields, prepared main-thread data, worker-resident backing,
worker scratch, handoff buffers, retained geometry and GPU resources. Some
buffers move between stages; some are shared by multiple consumers. Accounting
uses allocation identity so shared ownership is not mistaken for two physical
copies, and a transition is not mistaken for freeing memory that remains live.

The peak working set can be larger than the final picture:
\begin{equation}
 B_{\mathrm{peak}} \approx B_{\mathrm{native}}+B_{\mathrm{decoded}}
  +B_{\mathrm{workers}}+B_{\mathrm{geometry}}+B_{\mathrm{GPU}}
  +B_{\mathrm{runtime}}.
\end{equation}
This is an explanatory decomposition, not a direct measurement of whole-process
resident memory. Shared allocations, transfers and overlapping lifetimes must
be handled by the actual owner ledger. Runtime overhead includes browser and
native-library work that cannot be inferred from array sizes alone.

For scale, one $256\times256\times160$ float32 scalar contains 40 MiB of values.
Six such scalar buffers contain 240 MiB before coordinates, copies, validity,
scratch, geometry, graphics textures and application overhead. This arithmetic
is useful for planning; it is not a promise that a particular scene uses exactly
that amount or that a computer with that much free RAM can render it.

\begin{figure}[htbp]
\centering
\begin{tikzpicture}[
  box/.style={draw=NimbusBlue,rounded corners,fill=NimbusLight,align=center,
    text width=3.45cm,minimum height=1.1cm,font=\small},
  flow/.style={-{Latex[length=2mm]},thick,draw=NimbusBlue}]
\node[box] (native) {Native read reservation\\Source work and reply};
\node[box,right=0.7cm of native] (field) {Decoded field/cache\\Shared consumer ownership};
\node[box,right=0.7cm of field] (worker) {Preparation worker\\Backing and scratch};
\node[box,below=0.8cm of worker] (resident) {Admit atomically\\Resident geometry\\and textures};
\node[box,left=0.7cm of resident] (retire) {Retire or retain explicitly\\Release, transfer or cache};
\draw[flow] (native) -- (field);
\draw[flow] (field) -- (worker);
\draw[flow] (worker) -- (resident);
\draw[flow] (resident) -- (retire);
\draw[flow,dashed] (worker.south west) -- (retire.north east);
\end{tikzpicture}
\caption{The I/O preference governs the first reservation, not every box.
Preparation may publish only after its resident claims are accepted. Cancelled
work and deselected layers retire their claims; reusable shared data keeps an
explicit owner.}
\end{figure}


\subsection{Admission, replacement and retirement}
A preparation task reserves the resources it expects to need. On completion,
the owner must transition construction claims to resident claims successfully
before publishing the prepared scene. If resident admission fails, the new
scene must not become visible with unaccounted backing. This commit boundary
keeps resource state and displayed state consistent.

Deselecting a volume field, replacing a preparation request or unmounting a
plot retires its inactive bricks and worker resources. Unfinished superseded
workers are stopped when their work has no remaining owner. Completed reusable
worker backing can remain warm under an explicit cache owner; that is different
from retaining every field ever selected. The corresponding error state must
also identify the current failed request rather than preserve an obsolete
failure from another field.

These rules apply to single and nested volume preparation, and to independently
required background reads. A budget increase must reach each applicable owner;
retrying only the foreground while leaving its required background failed would
leave the feature broken. Shared data stays alive while another consumer
needs it, so a drop in one pane's usage need not make every memory counter zero.

\subsection{What the resource panel can and cannot certify}
The resource panel exposes tracked ownership, reservations, resident commitments
and high-water observations. The aggregate visualization-capacity policy remains
observe-only unless a matching qualification artifact supplies the required
whole-application, GPU-residency and opaque-worker headroom evidence. Separate
native I/O limits, queue bounds, texture limits and geometry/output limits still
apply.

Observe-only does not mean unlimited memory is safe. It means that the tracked
sum is not being advertised as a universally qualified whole-application cap.
Browser subprocesses, graphics drivers, native decoder scratch and transient
runtime allocations need measurement on the actual packaged application.
An observed small-scene success is not a safe upper bound for every dataset.

\subsection{ARM64, x86-64 and memory topology}
ARM64 and x86-64 describe instruction sets. Unified and discrete memory describe
where processors access storage. Neither distinction alone predicts Studio's
performance. An ARM64 application still needs the correct runtime libraries;
an x86-64 machine can have integrated graphics; host-visible GPU memory can
still have costs unlike a tightly unified memory system.

On a discrete GPU, texture upload can create a separate device allocation and
consume interconnect bandwidth. On a unified system, CPU and GPU work still need
proper synchronization and ownership, and concurrent access can compete for the
same physical bandwidth. JavaScript arrays, worker transfers, native decoder
buffers and graphics textures do not become one shared zero-copy allocation
merely because the hardware shares memory.

Measure the viewing computer separately from a remote compute node. A remote
GPU's memory does not raise Studio's local texture capacity. Compare source
identity, field set, sampling, visible layers, worker count and graphics backend
when measuring changes. Startup, repeated-frame preparation, steady rendering,
capture and file export are separate intervals with different bottlenecks.
